\documentclass[12pt]{amsart}
%\usepackage[margin=1in]{geometry}
\pagestyle{plain}

\usepackage{amsmath,amssymb}
\usepackage{enumerate}

\usepackage[shortlabels]{enumitem}
% Set spacing for all enumerate environments
\setlist[enumerate]{itemsep=6pt, topsep=6pt, parsep=0pt}

\usepackage{graphicx}

\usepackage[left=0.75in,right=1in,bottom=0.9in]{geometry}  
% \usepackage{mathptmx}      % use Times fonts if available on your TeX system
\usepackage[T1]{fontenc}

\usepackage[parfill]{parskip}
\renewcommand{\baselinestretch}{1}

% Macro definitions
\newcommand{\norm}[1]{\left\|#1\right\|}

%%%%%%%%%%%%%%%%%%%%%%%%%%%%%%%%%
% SECTION DIVIDERS
%%%%%%%%%%%%%%%%%%%%%%%%%%%%%%%%%
\newcommand{\divider}{
  \vspace{-1em} % Add space above
  \noindent
  \raisebox{-0.15ex}{\textbullet}%
  \hspace{0.5em}%
  \rule[0.5ex]{0.95\textwidth}{1.0pt}%
  \hspace{0.5em}%
  \raisebox{-0.15ex}{\textbullet}%
  \vspace{0em} % Adjust space below
}

\title{ESE 2030 QUIZZAM 3 : Spring 2026}
\author{prof-g \& meg}


\begin{document}
\maketitle

\begin{center}
    {\bf INSTRUCTIONS}
\end{center}

\begin{itemize}
\item 
No books, notes, or calculators. Use a writing utensil and logic. 

\item 
Please follow question instructions and fill in the bubble-sheet. 

\item 
Select one answer per problem.

\item 
Each problem has one best answer. 

\item 
In some problems, a wrong choice may be worth partial credit. 

\item 
Be careful to fill in the correct problem number. 

\item 
These questions are written carefully: if you detect an error, read closely and do your best.

\item 
If anything is unclear, use your best judgment. 

\item 
Cheating in any form will be dealt with severely. 
\end{itemize}

\newpage


% QUIZZAM 5 DRAFT
% TOTAL QUESTIONS: 10 (of target 20-25; Week 9 problems still needed)
% COVERAGE: WEEKS 7-9


{\bf PROBLEM 1:}
% WEEK = 7
% TOPICS = Matrix exponential, initial value problem, linearity of solutions, superposition
% NOTE: Rewritten stem for clarity (removed "particular" and "IVP" language)
%
Consider the system $\dfrac{d\mathbf{x}}{dt} = A\mathbf{x}$ where $A$ is a 
$3 \times 3$ matrix. Suppose you know two solutions:
\begin{itemize}
\item $\mathbf{x}_1(t)$ is the solution with initial condition $\mathbf{x}_1(0) = \mathbf{a}$
\item $\mathbf{x}_2(t)$ is the solution with initial condition $\mathbf{x}_2(0) = \mathbf{b}$
\end{itemize}
%
What is the solution $\mathbf{x}(t)$ with initial condition $\mathbf{x}(0) = 3\mathbf{a} - 2\mathbf{b}$?

\begin{enumerate}[(A)]
\item $3\mathbf{x}_1(t) - 2\mathbf{x}_2(t)$
\item $e^{At}(3\mathbf{a} - 2\mathbf{b})$, but this cannot be simplified 
      further without computing $e^{At}$
\item $(3-2t)\mathbf{x}_1(t) + t\,\mathbf{x}_2(t)$, by interpolation 
      between the two known solutions
\item $3\mathbf{x}_1(t) - 2\mathbf{x}_2(t)$, but only if $A$ is diagonalizable
\item Cannot be determined without knowing $A$ explicitly
\end{enumerate}

% \textbf{Correct Answer:} A
%
% \textbf{Explanation:} The solution operator $\mathbf{x}_0 \mapsto 
% e^{At}\mathbf{x}_0$ is \emph{linear} in the initial condition, because 
% matrix multiplication is linear. This means:
%   $e^{At}(3\mathbf{a} - 2\mathbf{b}) = 3e^{At}\mathbf{a} - 2e^{At}\mathbf{b} 
%     = 3\mathbf{x}_1(t) - 2\mathbf{x}_2(t)$
%
% No knowledge of $A$ is needed beyond the fact that the system is 
% linear. This is the superposition principle for linear ODEs: any 
% linear combination of solutions is again a solution, with the 
% corresponding linear combination of initial conditions.
%
% Note that this works for \emph{any} matrix $A$ --- diagonalizable or 
% not, real or complex eigenvalues, repeated eigenvalues or not. The 
% linearity comes from the structure of matrix multiplication, not from 
% any special property of $A$.
%
% \textbf{Partial credit:} B (correctly writes down the formal 
%   expression $e^{At}(3\mathbf{a} - 2\mathbf{b})$, which is true, 
%   but fails to recognize that linearity of $e^{At}$ allows this 
%   to be rewritten using the known solutions. The student understands 
%   the matrix exponential formula but misses the deeper point that 
%   known solutions can be combined without ever computing $e^{At}$ 
%   explicitly.)
%
% \textbf{Trick answer:} E (fails to recognize that linearity allows 
%   us to build new solutions from old ones; thinks every new initial 
%   condition requires recomputing from scratch)
%
% \textbf{Trick answer:} D (the superposition principle holds for all 
%   linear systems, not just diagonalizable ones. Linearity of the 
%   solution operator comes from linearity of matrix multiplication, 
%   which requires no assumptions about the eigenstructure of $A$. 
%   A student choosing this may conflate the techniques for 
%   \emph{computing} $e^{At}$ --- which do depend on eigenstructure 
%   --- with the structural fact that $e^{At}$ is linear.)
%
% \textbf{Trick answer:} C (invents a nonlinear interpolation scheme; 
%   the time-dependent coefficients $(3-2t)$ and $t$ would violate 
%   the ODE --- a linear combination of solutions with time-dependent 
%   coefficients is generally not a solution unless the coefficients 
%   are constant)
%
% \textbf{Trick answer:} B (see partial credit; formally correct but 
%   misses the simplification that makes the matrix exponential 
%   powerful in practice)

\divider

{\bf PROBLEM 2:}
% WEEK = 9
% TOPICS = Spectral radius, linear iteration, convergence to zero, eigenvalue magnitude vs sign
%
A $3 \times 3$ matrix $A$ has eigenvalues $\lambda_1 = 0.7$, 
$\lambda_2 = -0.9$, and $\lambda_3 = 0.2$.
%
Consider the iteration $\mathbf{x}_{k+1} = A\mathbf{x}_k$ starting 
from a generic initial condition $\mathbf{x}_0$.
%
What is the long-term behavior of $\mathbf{x}_k$?

\begin{enumerate}[(A)]
\item $\mathbf{x}_k$ may converge or diverge depending on the initial condition
\item $\mathbf{x}_k$ oscillates indefinitely without settling to a limit
\item $\mathbf{x}_k$ converges to a dominant eigenvector
\item $\mathbf{x}_k$ converges to $\mathbf{0}$ for every initial condition
\item $\mathbf{x}_k$ diverges due to oscillations at each step
\end{enumerate}

% \textbf{Correct Answer:} D
%
% \textbf{Explanation:} The key fact is that $A^k \to 0$ if and only if 
% every eigenvalue of $A$ satisfies $|\lambda_i| < 1$, i.e., the spectral 
% radius $\rho(A) < 1$. Here, $|\lambda_1| = 0.7$, $|\lambda_2| = 0.9$, 
% and $|\lambda_3| = 0.2$, so $\rho(A) = 0.9 < 1$.
%
% What happens mechanically: when $A$ is diagonalizable (or more generally), 
% the $k$-th iterate decomposes along eigenvectors as 
% $\mathbf{x}_k = c_1 \lambda_1^k \mathbf{v}_1 + c_2 \lambda_2^k 
% \mathbf{v}_2 + c_3 \lambda_3^k \mathbf{v}_3$. Each term decays because 
% $|\lambda_i|^k \to 0$. The term $(-0.9)^k$ does alternate in sign, but 
% its \emph{magnitude} $0.9^k \to 0$, so this component oscillates with 
% decreasing amplitude and still vanishes.
%
% This holds for \emph{every} initial condition, not just special ones: 
% every component of $\mathbf{x}_k$ along every eigendirection is 
% multiplied by a quantity with magnitude less than $1$ at each step.
%
% \textbf{Partial credit:} B (correctly identifies the sign-alternating 
%   behavior of $(-0.9)^k$, which is a real feature of the iteration: 
%   the component along $\mathbf{v}_2$ \emph{does} oscillate in sign 
%   at each step. The error is concluding that this oscillation persists 
%   forever at constant amplitude. In fact, $|{-0.9}|^k = 0.9^k \to 0$, 
%   so the oscillations are damped and die out. The student understands 
%   the qualitative effect of a negative eigenvalue but fails to connect 
%   magnitude $< 1$ with decay.)
%
% \textbf{Trick answer:} E (confuses the sign of $\lambda_2$ with the 
%   growth of iterates. A negative eigenvalue with $|\lambda| < 1$ 
%   produces \emph{decaying} sign-alternating behavior, not growing 
%   oscillations. The student conflates ``negative'' with ``unstable.'')
%
% \textbf{Trick answer:} C (two errors compounded. First, it identifies 
%   $\lambda_1 = 0.7$ as ``the largest eigenvalue'' --- but the dominant 
%   eigenvalue is determined by magnitude, and $|\lambda_2| = 0.9 > 0.7$. 
%   Second, it claims the iteration converges to a nonzero limit, which 
%   would require a dominant eigenvalue with $|\lambda| = 1$, not $< 1$. 
%   When all $|\lambda_i| < 1$, every component decays to zero; there is 
%   no surviving eigendirection.)
%
% \textbf{Trick answer:} B (see partial credit; the sign alternation is 
%   real but the student doesn't recognize damping)
%
% \textbf{Trick answer:} A (when $\rho(A) < 1$, convergence to 
%   $\mathbf{0}$ holds for \emph{all} initial conditions, not just 
%   some. The behavior is global, not initial-condition-dependent. 
%   Initial-condition dependence arises when some eigenvalues have 
%   $|\lambda| < 1$ and others have $|\lambda| > 1$, which is not 
%   the case here.)

\divider

{\bf PROBLEM 3:}
% WEEK = 8
% TOPICS = Real Jordan canonical form, block structure, complex conjugate 
%          2x2 blocks, Jordan chains, reading block boundaries
%
The $6 \times 6$ matrix ${\mathcal J}$ below is in real Jordan canonical form:
$${\mathcal J} = \begin{bmatrix} 
3 & 1 & 0 & 0 & 0 & 0 \\ 
0 & 3 & 0 & 0 & 0 & 0 \\ 
0 & 0 & 3 & 0 & 0 & 0 \\ 
0 & 0 & 0 & 0 & 0 & 0 \\ 
0 & 0 & 0 & 0 & -1 & -2 \\ 
0 & 0 & 0 & 0 & 2 & -1 
\end{bmatrix}$$
How many Jordan blocks does ${\mathcal J}$ have?
\begin{enumerate}[(A)]
\item $6$
\item $3$
\item $4$
\item $5$
\item This matrix is not in Jordan canonical form
\end{enumerate}
% \textbf{Correct Answer:} C
%
% \textbf{Explanation:} The four blocks are:
% \begin{enumerate}
%   \item A $2 \times 2$ Jordan block 
%     $\begin{bmatrix} 3 & 1 \\ 0 & 3 \end{bmatrix}$ for 
%     $\lambda = 3$ (the superdiagonal $1$ chains an eigenvector 
%     to a generalized eigenvector).
%   \item A $1 \times 1$ block $[3]$ for $\lambda = 3$.
%   \item A $1 \times 1$ block $[0]$ for $\lambda = 0$.
%   \item A $2 \times 2$ rotation block 
%     $\begin{bmatrix} -1 & -2 \\ 2 & -1 \end{bmatrix}$ encoding 
%     the complex conjugate pair $\lambda = -1 \pm 2i$. This has 
%     the standard form 
%     $\begin{bmatrix} \alpha & -\beta \\ \beta & \alpha 
%     \end{bmatrix}$ with $\alpha = -1$, $\beta = 2$.
% \end{enumerate}
% Total: $4$ blocks. The key subtleties are (i) the eigenvalue 
% $\lambda = 3$ produces \emph{two} separate blocks (a $2 \times 2$ 
% Jordan chain and a $1 \times 1$ block), not one combined block, 
% and (ii) the $2 \times 2$ rotation block is a \emph{single} block 
% in real Jordan form.
%
% \textbf{Partial credit:} None
%
% \textbf{Trick answer:} B --- Counts one block per distinct 
%   eigenvalue: one block for $\lambda = 3$, one for $\lambda = 0$, 
%   one for the complex pair. This conflates ``same eigenvalue'' 
%   with ``same block.'' In Jordan form, a single eigenvalue can 
%   produce multiple blocks --- here $\lambda = 3$ has a $2 \times 2$ 
%   Jordan block AND a separate $1 \times 1$ block (reflecting 
%   geometric multiplicity $2$ within algebraic multiplicity $3$).
%
% \textbf{Trick answer:} D --- Splits the $2 \times 2$ rotation 
%   block into two $1 \times 1$ blocks by reading $-1$ and $-1$ off 
%   its diagonal, yielding five blocks: the $2 \times 2$ Jordan 
%   block, $[3]$, $[0]$, $[-1]$, $[-1]$. This targets the student 
%   who does not recognize the rotation block as a single unit 
%   encoding a complex conjugate pair.
%
% \textbf{Trick answer:} A --- Treats every diagonal entry as its 
%   own $1 \times 1$ block, ignoring all block structure entirely. 
%   This is the ``I don't know what blocks are'' answer: six 
%   diagonal entries, six blocks.
%
% \textbf{Trick answer:} E --- Believes the off-diagonal entries 
%   $\pm 2$ in the lower-right $2 \times 2$ block violate Jordan 
%   form. This targets the student who only knows \emph{complex} 
%   Jordan form (where all blocks are upper triangular with 
%   eigenvalues on the diagonal) and does not recognize that 
%   \emph{real} Jordan form uses $2 \times 2$ rotation blocks to 
%   encode complex conjugate pairs using only real entries.


%%%%%%%%%%%%%%%%%%%%%%%%%%%%%%%%%%%%%%%%%%%%%%%%%%%%
\newpage
%%%%%%%%%%%%%%%%%%%%%%%%%%%%%%%%%%%%%%%%%%%%%%%%%%%%


{\bf PROBLEM 4:}
% WEEK = 9
% TOPICS = Graph Laplacian, symmetry, Spectral Theorem, kernel structure, 
%          row/column sums, why zero is always an eigenvalue
%
Let $G$ be an undirected graph on $n$ vertices with adjacency matrix $A$ 
and degree matrix $D = \operatorname{diag}(d_1, \ldots, d_n)$. The graph 
Laplacian is $L = D - A$.

Which of the following best explains why $\lambda = 0$ is always an 
eigenvalue of $L$, regardless of the graph?

\begin{enumerate}[(A)]
\item Every row of $L$ sums to zero, so the all-ones 
      vector $\mathbf{1}$ is in $\ker(L)$
\item $L = D - A$ is a difference of two matrices, and the subtraction 
      causes cancellation that forces a zero eigenvalue
\item The diagonal entries of $L$ are non-negative and the off-diagonal 
      entries are non-positive, which forces at least one eigenvalue 
      to be zero
\item $L$ is symmetric, and symmetric matrices always have at least one zero eigenvalue
\item None of the above
\end{enumerate}

% \textbf{Correct Answer:} A
%
% \textbf{Explanation:} The graph Laplacian $L = D - A$ has the 
% property that its rows sum to zero. This is because the diagonal 
% entry $L_{ii} = d_i$ (the degree of vertex $i$) and the off-diagonal 
% entries in row $i$ are $L_{ij} = -A_{ij}$, so:
%   $$\sum_j L_{ij} = d_i - \sum_{j \neq i} A_{ij} = d_i - d_i = 0$$
% since the degree of vertex $i$ equals the number of its neighbors.
%
% This means $L\mathbf{1} = \mathbf{0}$, where $\mathbf{1}$ is the 
% vector of all ones. By definition, $\mathbf{1}$ is an eigenvector 
% of $L$ with eigenvalue $0$. This argument works for any undirected 
% graph --- connected or disconnected, with any degree sequence.
%
% Moreover, since $L$ is symmetric (the graph is undirected, so 
% $A = A^T$, and $D$ is diagonal), the Spectral Theorem guarantees 
% that $L$ can be orthogonally diagonalized with all real eigenvalues. 
% It is also positive semi-definite (since 
% $\mathbf{x}^T L \mathbf{x} = \sum_{(i,j) \in E} (x_i - x_j)^2 
% \geq 0$), so all eigenvalues are $\geq 0$, and we have just shown 
% that $0$ is achieved.
%
% \textbf{Partial credit:} None
%
% \textbf{Trick answer:} D (false --- symmetric matrices do NOT 
%   necessarily have a zero eigenvalue. For example, the identity 
%   matrix $I$ is symmetric with all eigenvalues equal to $1$. 
%   Symmetry guarantees real eigenvalues and orthogonal eigenvectors, 
%   but says nothing about whether zero is among them. The zero 
%   eigenvalue of $L$ comes from the specific row-sum structure, 
%   not from symmetry alone.)
%
% \textbf{Trick answer:} B (the fact that $L$ is a ``difference'' 
%   does not explain the zero eigenvalue. Many differences of 
%   matrices are invertible --- e.g., $3I - I = 2I$ has no zero 
%   eigenvalue. The zero eigenvalue arises from the specific 
%   algebraic relationship between $D$ and $A$, not from the 
%   subtraction operation itself.)
%
% \textbf{Trick answer:} C (the sign pattern of $L$ does not force 
%   a zero eigenvalue. The matrix 
%   $\begin{bmatrix} 2 & -1 \\ -1 & 3 \end{bmatrix}$ has the 
%   same sign pattern as a Laplacian but is invertible, with 
%   eigenvalues $\frac{5 \pm \sqrt{5}}{2}$. The sign pattern 
%   describes the class of \emph{M-matrices}, not all of which 
%   are singular.)
%
% \textbf{Trick answer:} E 

\divider

{\bf PROBLEM 5:}
% WEEK = 7
% TOPICS = Matrix exponential, power series definition, nilpotent matrices, 
%          direct computation from definition
%
Which of the following is $e^{A}$ where 
$A = \begin{bmatrix} 0 & -2 & 0 \\ 0 & 0 & -2 \\ 0 & 0 & 0 \end{bmatrix}$?

\noindent
\begin{minipage}{0.33\textwidth}\centering
(A) $\begin{bmatrix} 1 & -2 & -2 \\ 0 & 1 & -2 \\ 0 & 0 & 1 \end{bmatrix}$
\end{minipage}%
\begin{minipage}{0.33\textwidth}\centering
(B) $\begin{bmatrix} 1 & -2 & 4 \\ 0 & 1 & -2 \\ 0 & 0 & 1 \end{bmatrix}$
\end{minipage}%
\begin{minipage}{0.33\textwidth}\centering
(C) $\begin{bmatrix} 1 & e^{-2} & 0 \\ 0 & 1 & e^{-2} \\ 0 & 0 & 1 \end{bmatrix}$
\end{minipage}

\medskip
\noindent
\begin{minipage}{0.5\textwidth}\centering
(D) $\begin{bmatrix} 1 & -2 & 2 \\ 0 & 1 & -2 \\ 0 & 0 & 1 \end{bmatrix}$
\end{minipage}%
\begin{minipage}{0.5\textwidth}\centering
(E) $\begin{bmatrix} 1 & -2 & 0 \\ 0 & 1 & -2 \\ 0 & 0 & 1 \end{bmatrix}$
\end{minipage}

% \textbf{Correct Answer:} D
%
% \textbf{Explanation:} The matrix $A$ is strictly upper triangular, 
% hence nilpotent. We compute:
%   $$A^2 = \begin{bmatrix} 0 & -2 & 0 \\ 0 & 0 & -2 \\ 0 & 0 & 0 
%           \end{bmatrix}
%           \begin{bmatrix} 0 & -2 & 0 \\ 0 & 0 & -2 \\ 0 & 0 & 0 
%           \end{bmatrix}
%         = \begin{bmatrix} 0 & 0 & 4 \\ 0 & 0 & 0 \\ 0 & 0 & 0 
%           \end{bmatrix}$$
% and $A^3 = A^2 \cdot A = 0$ (which can be verified directly, or 
% follows from the fact that an $n \times n$ strictly upper triangular 
% matrix is nilpotent of index at most $n$).
%
% Since $A^k = 0$ for all $k \geq 3$, the power series terminates:
%   $$e^A = I + A + \frac{1}{2!}A^2 
%         = \begin{bmatrix} 1 & 0 & 0 \\ 0 & 1 & 0 \\ 0 & 0 & 1 
%           \end{bmatrix}
%         + \begin{bmatrix} 0 & -2 & 0 \\ 0 & 0 & -2 \\ 0 & 0 & 0 
%           \end{bmatrix}
%         + \frac{1}{2}\begin{bmatrix} 0 & 0 & 4 \\ 0 & 0 & 0 \\ 
%           0 & 0 & 0 \end{bmatrix}
%         = \begin{bmatrix} 1 & -2 & 2 \\ 0 & 1 & -2 \\ 0 & 0 & 1 
%           \end{bmatrix}$$
%
% The problem tests the entire chain: recognizing nilpotency so the 
% series is finite, computing $A^2$ correctly (including the sign of 
% $(-2)(-2) = +4$), applying the $\frac{1}{k!}$ coefficients properly, 
% and assembling $I + A + \frac{1}{2}A^2$.
%
% \textbf{Partial credit:} NONE
%
% \textbf{Trick answer:} B (The student does the hard parts 
%   correctly --- nilpotent recognition, matrix squaring, series 
%   truncation --- but forgets the factorial denominators that 
%   distinguish $e^A$ from $\sum A^k$.)
%
% \textbf{Trick answer:} E (the student computes only $I + A$, 
%   believing the series terminates after the linear term. This 
%   means the student either thinks $A^2 = 0$ without checking 
%   --- confusing nilpotent index $2$ with index $3$ --- or 
%   doesn't realize that higher-order terms can contribute. The 
%   $(1,3)$ entry is $0$ instead of $2$, revealing that the 
%   $A^2$ contribution was missed entirely.)
%
% \textbf{Trick answer:} C 
% You do not simply exponentiate the terms
%
% \textbf{Trick answer:} A (the student correctly sets up 
%   $I + A + \frac{1}{2}A^2$ and applies the factorial, but 
%   makes a sign error in computing $A^2$: gets $(-2)(-2) = -4$ 
%   instead of $+4$ in the $(1,3)$ entry, yielding 
%   $\frac{1}{2}(-4) = -2$. This is a pure arithmetic error 
%   in matrix multiplication, where the student loses track of 
%   the double negative.)

\vspace{0.5in}

\divider



{\bf PROBLEM 6:}
% WEEK = 9
% TOPICS = Spectral Theorem, symmetric vs diagonalizable, real eigenvalues, 
%          orthogonal eigenvectors, guaranteed diagonalizability
%
A real $n \times n$ matrix $A$ is known to be symmetric.
%
Which of the following is a consequence of the Spectral Theorem that does 
\textbf{not} follow from diagonalizability alone?

\begin{enumerate}[(A)]
\item $\operatorname{tr}(A)$ equals the sum of the eigenvalues of $A$
\item $\det(A)$ equals the product of the eigenvalues of $A$
\item $A$ can be written as $V\Lambda V^{-1}$ for some invertible $V$ 
      and diagonal $\Lambda$
\item $A$ has $n$ eigenvalues with algebraic and geometric multiplicity one
\item The eigenvectors of $A$ can be chosen to form an orthonormal basis 
      for $\mathbb{R}^n$
\end{enumerate}

% \textbf{Correct Answer:} E
%
% \textbf{Explanation:} The Spectral Theorem for real symmetric matrices 
% states that $A = Q\Lambda Q^T$ where $Q$ is orthogonal and $\Lambda$ 
% is real diagonal. This means:
%   (1) All eigenvalues are real.
%   (2) There exists a \emph{full orthonormal} eigenbasis.
%   (3) $A$ is always diagonalizable, regardless of eigenvalue 
%       multiplicity.
%
% Among the choices, only (E) --- the existence of an orthonormal 
% eigenbasis --- requires symmetry. The others are either properties 
% of all matrices or of all diagonalizable matrices:
%
%   (A) Every $n \times n$ matrix (symmetric or not, diagonalizable 
%       or not) has $n$ eigenvalues counting multiplicity, because 
%       the characteristic polynomial always has degree $n$.
%
%   (B) This is the \emph{definition} of diagonalizability. It 
%       follows from diagonalizability alone; symmetry is not needed.
%
%   (C) $\operatorname{tr}(A) = \sum \lambda_i$ holds for \emph{every} 
%       square matrix, not just diagonalizable or symmetric ones. It 
%       follows from the relationship between the characteristic 
%       polynomial and the trace.
%
%   (D) $\det(A) = \prod \lambda_i$ also holds for every square 
%       matrix, for the same reason as (C).
%
% The conceptual point: diagonalizability gives you a basis of 
% eigenvectors; the Spectral Theorem upgrades this to an 
% \emph{orthonormal} basis. This is a strictly stronger condition. 
% A diagonalizable matrix with non-orthogonal eigenvectors (e.g., 
% $\begin{bmatrix} 2 & 1 \\ 0 & 3 \end{bmatrix}$) satisfies 
% (A)--(D) but not (E).
%
% \textbf{Partial credit:} None
%
% \textbf{Trick answer:} D (true for all matrices, not specific to 
%   symmetry or even to diagonalizability; students may think the 
%   Spectral Theorem is needed to guarantee $n$ eigenvalues exist)
%
% \textbf{Trick answer:} C (this IS diagonalizability --- precisely 
%   the property the question says is already assumed. A student 
%   who selects this fails to distinguish diagonalizability from 
%   the additional structure that symmetry provides.)
%
% \textbf{Trick answer:} A (true for all matrices; a student may 
%   associate the trace-eigenvalue relationship with diagonalization, 
%   but it holds universally)
%
% \textbf{Trick answer:} B (same reasoning as C; the 
%   determinant-eigenvalue relationship is universal)



%%%%%%%%%%%%%%%%%%%%%%%%%%%%%%%%%%%%%%%%%%%%%%%%%%%%
\newpage
%%%%%%%%%%%%%%%%%%%%%%%%%%%%%%%%%%%%%%%%%%%%%%%%%%%%


{\bf PROBLEM 7:}
% WEEK = 8
% TOPICS = Jordan canonical form, eigenvalues, algebraic/geometric multiplicity, block structure
%
A $4 \times 4$ matrix $A$ has only one eigenvalue, $\lambda = 3$, with 
algebraic multiplicity $4$. The eigenspace $\ker(A - 3I)$ has dimension $2$.
%
How many Jordan blocks appear in the Jordan form of $A$, and why?

\begin{enumerate}[(A)]
\item Two blocks: specifically one $3 \times 3$ block and one $1 \times 1$
      block, because the deficiency $4 - 2 = 2$ means the matrix $(A - 3I)$
      is not nilpotent of order $1$, forcing the larger block to have size $3$
\item Four $1 \times 1$ blocks, because the algebraic multiplicity is $4$
\item Two blocks, because the number of Jordan blocks for an eigenvalue 
      equals the dimension of its eigenspace
\item One $4 \times 4$ block, because there is only one eigenvalue
\item Cannot be determined without knowing the dimensions of $\ker(A - 3I)^k$ for $k \geq 2$
\end{enumerate}

% \textbf{Correct Answer:} C
%
% \textbf{Explanation:} A fundamental fact about Jordan form: the 
% number of Jordan blocks for a given eigenvalue equals the 
% \emph{geometric multiplicity} --- the dimension of the eigenspace. 
% Each Jordan block begins with an eigenvector, so the number of 
% independent eigenvectors is exactly the number of blocks.
%
% Here, $\dim(\ker(A - 3I)) = 2$, so there are exactly 2 Jordan 
% blocks for $\lambda = 3$. Since $A$ is $4 \times 4$ and has only 
% one eigenvalue, these two blocks must account for all 4 dimensions. 
% The blocks could be $2 \times 2$ and $2 \times 2$, or $3 \times 3$ 
% and $1 \times 1$ --- this cannot be determined without further 
% information (specifically, $\dim(\ker(A - 3I)^2)$). But the 
% \emph{number} of blocks is determined: it is $2$.
%
% \textbf{Partial credit:} A (correctly identifies that there are 
%   two blocks, which is the key insight --- the block count equals 
%   the geometric multiplicity. The partition $3 \times 3$ plus 
%   $1 \times 1$ is also a genuinely valid possibility. The only 
%   error is overcommitting to this specific partition when two 
%   $2 \times 2$ blocks is equally consistent with the given data. 
%   The student understands Jordan structure well but goes one step 
%   too far without the additional information needed to pin down 
%   the sizes. Choice (B) correctly states the count without 
%   overspecifying.)
%
% \textbf{Trick answer:} D (confuses ``one eigenvalue'' with ``one 
%   block.'' Having a single eigenvalue constrains all blocks to 
%   share that eigenvalue, but doesn't mean there's only one block. 
%   The number of blocks depends on the eigenspace dimension, not 
%   the number of distinct eigenvalues.)
%
% \textbf{Trick answer:} B (this would mean $A$ is diagonalizable 
%   --- four $1 \times 1$ blocks is the same as a diagonal matrix. 
%   But if the eigenspace is only $2$-dimensional while the algebraic 
%   multiplicity is $4$, the matrix is \emph{not} diagonalizable, 
%   so blocks of size $> 1$ must appear.)
%
% \textbf{Trick answer:} E (the block count \emph{is} determined by 
%   the eigenspace dimension alone --- it equals the geometric 
%   multiplicity. While it's true that generalized eigenspaces are 
%   needed to determine block \emph{sizes}, the question only asks 
%   about the number of blocks. A student selecting (E) fails to 
%   recognize what the geometric multiplicity directly tells us.)

\divider

{\bf PROBLEM 8:}
% WEEK = 9
% TOPICS = Stochastic matrices, stationary distribution, convergence, 
%          Perron-Frobenius, dominant eigenvalue
%
Consider the $3 \times 3$ positive stochastic matrix with stationary distribution:
$$P = \begin{bmatrix} 
0.4 & 0.2 & 0.1 \\ 
0.4 & 0.5 & 0.3 \\ 
0.2 & 0.3 & 0.6 
\end{bmatrix} \quad \Rightarrow \quad 
\boldsymbol{\pi} = \begin{pmatrix} 0.2 \\ 0.4 \\ 0.4 \end{pmatrix}$$
%
If %the system starts in state $1$ with certainty, so that 
$\mathbf{x}_0 = \begin{pmatrix} 1 \\ 0 \\ 0 \end{pmatrix}$, 
what is $\displaystyle\lim_{k \to \infty} P^k \mathbf{x}_0$?
\begin{enumerate}[(A)]
\item $\begin{pmatrix} 0.2 \\ 0.4 \\ 0.4 \end{pmatrix}$
\item The limit depends on $\mathbf{x}_0$; knowing 
      $\boldsymbol{\pi}$ alone is not enough to determine it
\item $\begin{pmatrix} 1/3 \\ 1/3 \\ 1/3 \end{pmatrix}$
\item $\begin{pmatrix} 0.4 \\ 0.4 \\ 0.2 \end{pmatrix}$
\item The limit does not necessarily exist
\end{enumerate}
% \textbf{Correct Answer:} A
%
% \textbf{Explanation:} By Perron-Frobenius, a positive stochastic 
% matrix has a unique stationary distribution $\boldsymbol{\pi}$, 
% which is the eigenvector for the dominant eigenvalue $\lambda = 1$. 
% All other eigenvalues satisfy $|\lambda_i| < 1$, so the 
% components of $\mathbf{x}_0$ along non-dominant eigenvectors 
% decay geometrically:
%   $$P^k\mathbf{x}_0 = c_1 \boldsymbol{\pi} + 
%     c_2 \lambda_2^k \mathbf{v}_2 + c_3 \lambda_3^k \mathbf{v}_3 
%     \;\;\longrightarrow\;\; c_1 \boldsymbol{\pi}$$
% Since $\mathbf{x}_0$ is a probability distribution (entries sum 
% to $1$), the coefficient $c_1 = 1$, giving 
% $\lim_{k\to\infty} P^k\mathbf{x}_0 = \boldsymbol{\pi}$.
% This holds for \emph{any} initial probability distribution.
%
% \textbf{Partial credit:} None
%
% \textbf{Trick answer:} D --- This is $P\mathbf{x}_0$, the first 
%   column of $P$: the distribution after \emph{one} step, not 
%   infinitely many. Particularly deceptive here because the entries 
%   $0.4, 0.4, 0.2$ are the same three numbers as 
%   $\boldsymbol{\pi} = (0.2, 0.4, 0.4)^T$ rearranged. A student 
%   who computes one step and sees familiar-looking numbers may 
%   convince themselves they have the answer without checking the 
%   order.
%
% \textbf{Trick answer:} C --- Confuses ``equilibrium'' with 
%   ``uniform.'' The misconception is that long-run behavior must 
%   treat all states equally. The stationary distribution reflects 
%   the structure of $P$; uniform is stationary only for doubly 
%   stochastic matrices (both rows \emph{and} columns sum to $1$).
%
% \textbf{Trick answer:} B --- Sounds sophisticated: ``the limit 
%   depends on the initial condition.'' This is true for dynamical 
%   systems with multiple attractors (e.g., when $|\lambda_1| = 
%   |\lambda_2|$), but for a positive stochastic matrix, 
%   Perron-Frobenius guarantees a \emph{unique} stationary 
%   distribution that attracts \emph{every} probability 
%   distribution. The initial state affects the transient behavior 
%   (how fast convergence occurs) but not the limiting state. 
%   A student who selects this has some dynamical systems intuition 
%   but does not know the convergence theorem.
%
% \textbf{Trick answer:} E --- Reverses the logic of stationarity. 
%   The equation $P\boldsymbol{\pi} = \boldsymbol{\pi}$ says that 
%   $\boldsymbol{\pi}$ is a \emph{fixed point}: if you start there, 
%   you stay. The convergence theorem says something much stronger: 
%   you arrive at $\boldsymbol{\pi}$ from \emph{any} starting 
%   distribution. A student selecting (E) understands stationarity 
%   as a local property but misses the global attractivity.


%%%%%%%%%%%%%%%%%%%%%%%%%%%%%%%%%%%%%%%%%%%%%%%%%%%%
\newpage
%%%%%%%%%%%%%%%%%%%%%%%%%%%%%%%%%%%%%%%%%%%%%%%%%%%%


{\bf PROBLEM 9:}
% WEEK = 7
% TOPICS = Eigenvectors, linear independence, eigenvalues, scalar multiples, sums of eigenvectors
% NOTE: Kept as an easier problem for the quiz.
%
Let $A$ be a $3 \times 3$ matrix with eigenvalues $\lambda_1 = 5$, 
$\lambda_2 = -2$, and $\lambda_3 = 5$ (note: $\lambda_1 = \lambda_3$). 
Let $\mathbf{v}_1$ be an eigenvector for $\lambda_1 = 5$ and $\mathbf{v}_2$ 
an eigenvector for $\lambda_2 = -2$.
%
Which of the following \textbf{must} be true?
\vspace{-0.1in}
\begin{enumerate}[(A)]
\item $\mathbf{v}_1$ and $\mathbf{v}_2$ are orthogonal
\item $\mathbf{v}_1 + \mathbf{v}_2$ is an eigenvector of $A$ with 
      eigenvalue $5 + (-2) = 3$
\item $\mathbf{v}_1$ and $\mathbf{v}_2$ are linearly independent
\item $A$ must be diagonalizable, since it has three eigenvalues 
      counting multiplicity
\item Any nonzero scalar multiple $c\mathbf{v}_1$ is also an eigenvector 
      of $A$, but with a eigenvalue $5c$
\end{enumerate}

% \textbf{Correct Answer:} C
%
% \textbf{Explanation:} Eigenvectors corresponding to \emph{distinct} 
% eigenvalues are always linearly independent. Since $\lambda_1 = 5 
% \neq -2 = \lambda_2$, the eigenvectors $\mathbf{v}_1$ and 
% $\mathbf{v}_2$ must be linearly independent regardless of any 
% other properties of $A$. This is a fundamental theorem about 
% eigenspaces: eigenspaces for distinct eigenvalues intersect only 
% at $\mathbf{0}$.
%
% This fact is the structural foundation that makes diagonalization 
% possible when a matrix has $n$ distinct eigenvalues: the $n$ 
% eigenvectors are automatically linearly independent, forming a basis.
%
% \textbf{Partial credit:} None
%
% \textbf{Trick answer:} A (orthogonality 
%   requires symmetry, which is not given)
%
% \textbf{Trick answer:} B (eigenvalues do not ``add'' when 
%   eigenvectors are summed. In fact, $\mathbf{v}_1 + \mathbf{v}_2$ 
%   is generally \emph{not} an eigenvector at all when $\lambda_1 
%   \neq \lambda_2$: we have $A(\mathbf{v}_1 + \mathbf{v}_2) = 
%   5\mathbf{v}_1 + (-2)\mathbf{v}_2$, which is not a scalar 
%   multiple of $\mathbf{v}_1 + \mathbf{v}_2$ unless $5 = -2$. 
%   This is a common misconception that treats eigenvalue computations 
%   as ``distributing'' over vector addition.)
%
% \textbf{Trick answer:} D (having three eigenvalues counting 
%   multiplicity is necessary but not sufficient for 
%   diagonalizability. The repeated eigenvalue $\lambda = 5$ has 
%   algebraic multiplicity 2, and $A$ is diagonalizable only if 
%   the geometric multiplicity also equals 2 --- i.e., the eigenspace 
%   for $\lambda = 5$ is two-dimensional. Without this information, 
%   diagonalizability cannot be concluded.)
%
% \textbf{Trick answer:} E (scalar multiples of eigenvectors are 
%   indeed eigenvectors, but the eigenvalue does \emph{not} change: 
%   $A(c\mathbf{v}_1) = cA\mathbf{v}_1 = c(5\mathbf{v}_1) = 
%   5(c\mathbf{v}_1)$. The eigenvalue remains $5$, not $5c$. 
%   This tests whether students understand that the eigenvalue is 
%   a property of the \emph{direction}, not the magnitude of the 
%   eigenvector.)

\divider


{\bf PROBLEM 10:}
% WEEK = 7
% TOPICS = Diagonalization, eigenvalues, geometric vs algebraic multiplicity, eigenbasis, geometric interpretation
%
A $3 \times 3$ matrix $A$ has eigenvalues $\lambda_1 = 3$, $\lambda_2 = 3$, and $\lambda_3 = -1$. 
{\bf Fact:} $A$ is diagonalizable.
What does diagonalizability of $A$ tell us that the eigenvalues alone do not?

\vspace{-0.1in}
\begin{enumerate}[(A)]
\item That $\operatorname{tr}(A) = 5$
\item That $A$ is symmetric, since only symmetric matrices can have repeated eigenvalues and still be diagonalizable
\item That there exists a basis of $\mathbb{R}^3$ consisting entirely of eigenvectors of $A$
\item That $\det(A) = -9$
\item That $A$ is invertible, since $e^A$ must be well-defined and nonsingular
\end{enumerate}

% \textbf{Correct Answer:} C
%
% \textbf{Explanation:} The eigenvalues alone --- even knowing that $\lambda = 3$ 
% is repeated --- do not tell us whether $A$ is diagonalizable. 
% Diagonalizability requires that the eigenspace for $\lambda = 3$ be 
% two-dimensional (geometric multiplicity = algebraic multiplicity). 
% When this holds, there exist two linearly independent eigenvectors 
% for $\lambda = 3$ alongside one eigenvector for $\lambda = -1$, 
% forming a full eigenbasis for $\mathbb{R}^3$. This is the defining 
% content of diagonalizability: $\mathbb{R}^3$ has a basis in which 
% $A$ acts by pure scaling on each basis vector.
%
% In contrast, the other choices are all deducible from the eigenvalues 
% alone, without knowing whether $A$ is diagonalizable:
%   - (A): $\det(A) = 3 \cdot 3 \cdot (-1) = -9 \neq 0$, so invertibility 
%     follows from the eigenvalues.
%   - (B): $\operatorname{tr}(A) = 3 + 3 + (-1) = 5$ from the eigenvalues.
%   - (D): $\det(A) = 3 \cdot 3 \cdot (-1) = -9$ from the eigenvalues.
%   - (E): Symmetry is a stronger condition than diagonalizability; a matrix 
%     can be diagonalizable without being symmetric (e.g., any matrix with 
%     distinct eigenvalues is diagonalizable but need not be symmetric).
%
% The question targets the conceptual distinction between information 
% encoded in the eigenvalues (trace, determinant, invertibility, 
% characteristic polynomial) versus information that requires knowing 
% about the eigenspaces (diagonalizability, existence of an eigenbasis).
%
% \textbf{Partial credit:} None
%
% \textbf{Trick answer:} E (true statement, but follows from the eigenvalues 
%   alone --- $\det \neq 0$ since no eigenvalue is zero. A student who 
%   selects this doesn't distinguish eigenvalue-level information from 
%   eigenspace-level information.)
%
% \textbf{Trick answer:} A (true statement, but trace equals the sum of 
%   eigenvalues regardless of diagonalizability. Same confusion as (A).)
%
% \textbf{Trick answer:} D (true statement, but determinant equals the 
%   product of eigenvalues regardless of diagonalizability. Same confusion 
%   as (A) and (B).)
%
% \textbf{Trick answer:} B (tempting because symmetric matrices are always 
%   diagonalizable, so students may conflate the two. But the implication 
%   only goes one way: symmetric $\Rightarrow$ diagonalizable, not the 
%   converse. Diagonalizability is strictly weaker than symmetry.)

\divider

{\bf PROBLEM 11:}
% WEEK = 9
% TOPICS = Stochastic matrices, eigenvalue properties, column sums, 
%          spectral radius bound, common misconceptions
%
A $3 \times 3$ matrix $P$ has non-negative entries and every column 
sums to $1$.
%
Which of the following \textbf{must} be true?
\vspace{-0.1in}
\begin{enumerate}[(A)]
\item All eigenvalues of $P$ are real
\item The columns of $P$ are linearly independent
\item Every eigenvalue of $P$ satisfies $0 \leq \lambda \leq 1$
\item $\det(P) = 1$
\item None of the above
\end{enumerate}

% \textbf{Correct Answer:} E
%
% \textbf{Explanation:} Each of (A)--(D) fails for some valid stochastic 
% matrix. The properties that DO always hold for a stochastic matrix 
% are: $\lambda = 1$ is always an eigenvalue, and every eigenvalue 
% satisfies $|\lambda| \leq 1$ (a bound on \emph{magnitude}, not on 
% the eigenvalue itself). None of the four listed statements captures 
% these correctly.
%
%   (A) FALSE: Stochastic matrices can have complex eigenvalues. The 
%       cyclic permutation matrix 
%       $\begin{bmatrix} 0 & 0 & 1 \\ 1 & 0 & 0 \\ 0 & 1 & 0 
%       \end{bmatrix}$ is stochastic (non-negative entries, columns 
%       sum to $1$) with eigenvalues $1, \, e^{2\pi i/3}, \, 
%       e^{-2\pi i/3}$ --- a conjugate pair of non-real complex 
%       numbers. Symmetry would guarantee real eigenvalues, but 
%       stochastic matrices are generally not symmetric.
%
%   (B) FALSE: This is the most instructive wrong answer. The correct 
%       bound on eigenvalues of a stochastic matrix is $|\lambda| \leq 1$ 
%       --- a constraint on \emph{magnitude}. This allows negative real 
%       eigenvalues (e.g., $\lambda = -0.5$) and complex eigenvalues 
%       (e.g., $\lambda = e^{2\pi i/3}$), as long as the modulus does 
%       not exceed $1$. The statement $0 \leq \lambda \leq 1$ incorrectly 
%       restricts eigenvalues to the interval $[0,1]$ on the real line, 
%       which excludes both negative and complex eigenvalues that are 
%       perfectly valid.
%
%   (C) FALSE: $\det(P)$ equals the product of eigenvalues, and there 
%       is no reason this product must be $1$. For instance, a stochastic 
%       matrix with eigenvalues $1, 0.5, -0.3$ has $\det(P) = -0.15$. 
%       A stochastic matrix with a zero eigenvalue has $\det(P) = 0$.
%
%   (D) FALSE: Stochastic matrices can be singular. The matrix 
%       $\frac{1}{2}\begin{bmatrix} 1 & 1 & 1 \\ 1 & 1 & 1 \\ 
%       0 & 0 & 0 \end{bmatrix}$ is not valid (third row sums are 
%       wrong), but the matrix with two identical columns, such as 
%       $P = \begin{bmatrix} 0.5 & 0.5 & 0.3 \\ 0.3 & 0.3 & 0.4 \\ 
%       0.2 & 0.2 & 0.3 \end{bmatrix}$, is stochastic with linearly 
%       dependent columns (columns $1$ and $2$ are equal), so 
%       $\det(P) = 0$.
%
% \textbf{Partial credit:} C (the student knows there is a bound 
%   involving $1$ on the eigenvalues of stochastic matrices, and 
%   correctly identifies $1$ as the upper bound. The error is writing 
%   the constraint as $0 \leq \lambda \leq 1$ rather than $|\lambda| 
%   \leq 1$. This is a genuine partial understanding: the student 
%   recalls the spectral radius bound $\rho(P) \leq 1$ but 
%   misinterprets it as a pointwise inequality on the real line 
%   rather than a constraint on magnitude. This reflects confusion 
%   between ``eigenvalue $\leq 1$'' and ``eigenvalue magnitude 
%   $\leq 1$'' --- a subtle but important distinction.)
%
% \textbf{Trick answer:} A (real eigenvalues hold for symmetric 
%   matrices by the Spectral Theorem, but stochastic matrices are 
%   not generally symmetric. A student selecting this may be 
%   conflating the two classes.)
%
% \textbf{Trick answer:} C (see partial credit; the most common 
%   error --- knowing the right bound but stating it incorrectly)
%
% \textbf{Trick answer:} D (a student may reason: ``columns sum to 
%   $1$, so the determinant should be $1$.'' This confuses the 
%   column-sum constraint with a determinant constraint. The column 
%   sums relate to the left eigenvector $\mathbf{1}^T$ for 
%   $\lambda = 1$, not to the product of all eigenvalues.)
%
% \textbf{Trick answer:} B (a student may assume that a ``well-behaved'' 
%   probability transition matrix should be invertible. In fact, 
%   stochastic matrices with repeated columns --- representing states 
%   with identical transition behavior --- are singular and arise 
%   naturally in applications.)

\divider

{\bf PROBLEM 12:}
% WEEK = 9
% TOPICS = Perron-Frobenius theorem, positive matrices, uniqueness of 
%          positive eigenvector, stochastic matrices, Markov chain convergence
%
A positive stochastic matrix $P$ (all entries strictly positive, columns 
sum to $1$) has eigenvalues $\lambda_i$.
The Perron-Frobenius theorem guarantees several properties of $P$.
Which of the following is \textbf{not} among them?
\vspace{-0.1in}
\begin{enumerate}[(A)]
\item Every other eigenvalue satisfies $|\lambda_i| < 1$
\item The eigenvalue $\lambda_1 = 1$ is simple (algebraic multiplicity 
      one)
\item The eigenvalue $\lambda_1 = 1$ has a corresponding eigenvector 
      with all entries strictly positive
\item For any initial probability distribution $\mathbf{x}_0$, the 
      iterates $P^k\mathbf{x}_0$ converge to the unique stationary 
      distribution as $k \to \infty$
\item None of the above : all these properties are guaranteed.
\end{enumerate}

% \textbf{Correct Answer:} E
%
% \textbf{Explanation:} Perron-Frobenius guarantees a great deal 
% about the \emph{dominant} eigenvalue and its eigenvector, but 
% its guarantees about the remaining eigenvectors are much more 
% limited. Specifically, for a positive matrix:
%
%   (A) TRUE: The Perron eigenvector (for $\lambda_1 = \rho(P)$) 
%       can be chosen with all entries strictly positive. This is 
%       a central conclusion of the theorem.
%
%   (B) TRUE: For a positive (not merely non-negative) matrix, 
%       the Perron eigenvalue is simple --- it has algebraic 
%       multiplicity one. This is the ``uniqueness'' part of 
%       Perron-Frobenius.
%
%   (C) TRUE: For a positive stochastic matrix, the Perron eigenvalue 
%       is $1$, and all other eigenvalues are strictly less than $1$ 
%       in magnitude. This is what makes convergence possible.
%
%   (D) TRUE: Since $|\lambda_i| < 1$ for $i \geq 2$, the components 
%       of $P^k\mathbf{x}_0$ along non-dominant eigenvectors decay 
%       to zero, leaving only the stationary distribution. The 
%       positivity of the Perron eigenvector combined with $\lambda_1 = 1$ 
%       and simplicity ensures convergence to a unique limit.
%
%
% \textbf{Partial credit:} None
%
% \textbf{Trick answer:} C (this IS guaranteed by Perron-Frobenius 
%   and is perhaps the theorem's most famous conclusion. A student 
%   selecting this may doubt that an eigenvector can have a sign 
%   constraint, but positivity of the Perron eigenvector is the 
%   theorem's core content.)
%
% \textbf{Trick answer:} B (simplicity of the Perron eigenvalue is 
%   guaranteed for positive matrices. A student might confuse this 
%   with the weaker result for non-negative matrices, where the 
%   Perron eigenvalue can have higher multiplicity in reducible 
%   cases. But the problem specifies a positive matrix.)
%
% \textbf{Trick answer:} A (this follows from the combination of 
%   Perron-Frobenius and the stochastic property. A student might 
%   doubt this if they recall that non-negative matrices can have 
%   multiple eigenvalues on the spectral circle, but strict 
%   positivity rules this out.)
%
% \textbf{Trick answer:} D (this is the convergence theorem for 
%   positive stochastic matrices. A student might think initial 
%   conditions matter for the limit, but the uniqueness of the 
%   stationary distribution and $|\lambda_i| < 1$ for $i \geq 2$ 
%   ensure universal convergence.)


%%%%%%%%%%%%%%%%%%%%%%%%%%%%%%%%%%%%%%%%%%%%%%%%%%%%
\newpage
%%%%%%%%%%%%%%%%%%%%%%%%%%%%%%%%%%%%%%%%%%%%%%%%%%%%


{\bf PROBLEM 13:}
% WEEK = 9
% TOPICS = Spectral Theorem (unstated), symmetric matrices, orthogonal 
%          eigenbasis, eigenbasis coordinate computation, dot products vs inversion
%
Let $A$ be a real symmetric $n \times n$ matrix, and let $B$ be a 
general (non-symmetric) diagonalizable $n \times n$ matrix. Both 
have $n$ linearly independent eigenvectors.
To express a vector $\mathbf{x} \in \mathbb{R}^n$ in the eigenbasis 
of each matrix, what key simplification does the symmetry of $A$ provide?
\begin{enumerate}[(A)]
\item The eigenvalues of $A$ are guaranteed real, so the eigenbasis 
      coordinates of $\mathbf{x}$ are real-valued --- unlike the 
      general case where coordinates may be complex
\item Both $A$ and $B$ require computing $V^{-1}\mathbf{x}$ to find 
      eigenbasis coordinates; symmetry does not simplify this step
\item The diagonalization of $A$ is unique, whereas 
      $B = V\Lambda V^{-1}$ has infinitely many choices of $V$
\item The eigenvectors of $A$ can always be chosen orthonormal, so 
      the eigenbasis coordinates are $c_i = \mathbf{q}_i^T\mathbf{x}$ 
      --- no matrix inversion is required
\item Symmetric matrices always have distinct eigenvalues, so there 
      is exactly one eigenbasis up to scaling
\end{enumerate}
% \textbf{Correct Answer:} D
%
% \textbf{Explanation:} The Spectral Theorem guarantees that a real 
% symmetric matrix possesses an orthonormal eigenbasis 
% $\{\mathbf{q}_1, \ldots, \mathbf{q}_n\}$. Forming 
% $Q = [\mathbf{q}_1 \cdots \mathbf{q}_n]$, we get $A = Q\Lambda Q^T$ 
% with $Q^{-1} = Q^T$. To find the eigenbasis coordinates of 
% $\mathbf{x}$:
%   $$\mathbf{c} = Q^T\mathbf{x} \quad \Longrightarrow \quad 
%     c_i = \mathbf{q}_i^T\mathbf{x}$$
% Each coordinate is a dot product --- an $O(n)$ operation requiring 
% no matrix inversion. For the general case $B = V\Lambda V^{-1}$, 
% the coordinates are $\mathbf{c} = V^{-1}\mathbf{x}$, which requires 
% solving a linear system ($O(n^2)$ or $O(n^3)$) and may be 
% ill-conditioned.
%
% The student must supply the unstated intermediate step: 
% symmetric $\Rightarrow$ Spectral Theorem $\Rightarrow$ orthonormal 
% eigenbasis $\Rightarrow$ $Q^{-1} = Q^T$ $\Rightarrow$ coordinates 
% are dot products.
%
% \textbf{Partial credit:} A (small --- the eigenvalues of a real 
%   symmetric matrix ARE guaranteed real, so this is a true statement; 
%   however, it does not address the computational simplification for 
%   finding eigenbasis coordinates. A general real matrix can also have 
%   all real eigenvalues and still require $V^{-1}$ for coordinate 
%   computation. The student recognizes a property of symmetric 
%   matrices but picks the wrong one.)
%
% \textbf{Trick answer:} C (the orthogonal diagonalization is NOT 
%   unique: within each eigenspace, any orthonormal basis may be 
%   chosen. The general diagonalization is also non-unique for the 
%   same reason. Uniqueness is not the distinguishing advantage.)
%
% \textbf{Trick answer:} E (false --- symmetric matrices can have 
%   repeated eigenvalues. The Spectral Theorem still guarantees a 
%   full orthonormal eigenbasis even with repetition, but distinctness 
%   is not guaranteed.)
%
% \textbf{Trick answer:} B (false --- this is precisely what symmetry 
%   DOES simplify. A student who does not know the Spectral Theorem 
%   may reason that both cases are ``just diagonalizable'' and conclude 
%   there is no difference. This distractor targets that gap.)

\divider

{\bf PROBLEM 14:}
% WEEK = 9 (with synthesis from Week 7)
% TOPICS = Dominant eigenvalue, linear recurrence as matrix iteration, 
%          convergence of ratios, power method, Fibonacci sequence
%
The Fibonacci sequence $F_0 = 0, \; F_1 = 1, \; F_{k+1} = F_k + F_{k-1}$ 
can be written in matrix form:
$$\begin{pmatrix} F_{k+1} \\ F_k \end{pmatrix} 
  = \underbrace{\begin{bmatrix} 1 & 1 \\ 1 & 0 \end{bmatrix}}_{A}
    \begin{pmatrix} F_k \\ F_{k-1} \end{pmatrix}
\quad \Rightarrow \quad \lambda_1 = \dfrac{1+\sqrt{5}}{2} \approx 1.618...\,\, ; \,\, \lambda_2= \dfrac{1-\sqrt{5}}{2} \approx -0.618...    
$$
It is a classical fact that $\dfrac{F_{k+1}}{F_k} \to \lambda_1 \approx 1.618\ldots$ as 
$k \to \infty$.
%
Which of the following best explains \textbf{why}?

\begin{enumerate}[(A)]
\item $\lambda_1$ is the only positive eigenvalue of $A$, and Fibonacci 
      numbers are positive, so the ratio must converge to the positive 
      eigenvalue
\item The trace of $A$ is $1$, so the sum of eigenvalues is $1$; as 
      $k$ grows, this forces the ratio $F_{k+1}/F_k$ toward the larger 
      eigenvalue
\item $A$ is symmetric, so the Spectral Theorem guarantees that the 
      dominant eigenvalue controls the long-term ratio
\item $\lambda_1$ satisfies $\lambda_1^2 = \lambda_1 + 1$, which is the same 
      relation as $F_{k+1} = F_k + F_{k-1}$ after dividing by $F_k$; 
      therefore any sequence obeying this recurrence must converge to 
      $\lambda_1$
\item Since $|\lambda_2| < 1 < |\lambda_1|$, the component of 
      $\mathbf{x}_k = A^k\mathbf{x}_0$ along the other 
      eigenvector decays relative to the component along $\lambda_1$'s 
      eigenvector, so $\mathbf{x}_k$ becomes approximately parallel to 
      $\lambda_1$'s eigenvector, and multiplication by $A$ then scales by 
      approximately $\lambda_1$
\end{enumerate}

% \textbf{Correct Answer:} E
%
% \textbf{Explanation:} Writing the state vector in the eigenbasis, 
% $\mathbf{x}_0 = c_1\mathbf{v}_1 + c_2\mathbf{v}_2$ where 
% $\mathbf{v}_1, \mathbf{v}_2$ are eigenvectors for $\lambda_1$ and 
% $\lambda_2$ respectively. Then:
%   $$\mathbf{x}_k = A^k\mathbf{x}_0 
%     = c_1\lambda_1^k\,\mathbf{v}_1 + c_2(\lambda_2)^k\,\mathbf{v}_2$$
%
% Since $|\lambda_2| \approx 0.618 < 1$ and $|\lambda_1| \approx 
% 1.618 > 1$, the first term grows while the second decays. For 
% large $k$:
%   $$\mathbf{x}_k \approx c_1\lambda_1^k\,\mathbf{v}_1$$
%
% Once $\mathbf{x}_k$ is approximately a multiple of $\mathbf{v}_1$, 
% one more multiplication by $A$ scales it by $\lambda_1$, meaning 
% both entries --- $F_{k+1}$ and $F_k$ --- grow by the same factor 
% $\lambda_1$. Their ratio therefore approaches $\lambda_1$.
%
% This is exactly the mechanism behind the power method: repeated 
% multiplication by a matrix causes the iterate to align with the 
% dominant eigenvector, and the ratio of successive iterates converges 
% to the dominant eigenvalue. The Fibonacci sequence is a natural 
% instance of this phenomenon.
%
% \textbf{Partial credit:} NONE
%
% \textbf{Trick answer:} D  (contains a true algebraic observation: 
%   dividing $F_{k+1} = F_k + F_{k-1}$ by $F_k$ gives 
%   $F_{k+1}/F_k = 1 + F_{k-1}/F_k$, and if the ratio converges 
%   to some limit $L$, then $L = 1 + 1/L$, i.e., $L^2 = L + 1$, 
%   which $\lambda_1$ satisfies. This argument correctly identifies 
%   the limit \emph{assuming convergence}, but does not explain 
%   \emph{why} convergence occurs. The claim ``any sequence obeying 
%   this recurrence must converge'' is false without the eigenvalue 
%   analysis: one must verify that $c_1 \neq 0$ and that the 
%   subdominant component actually decays. The fixed-point equation 
%   alone does not rule out oscillation, divergence, or convergence 
%   to $\hat{\lambda_1}$. A student selecting this has a correct 
%   algebraic intuition but has not identified the mechanism --- 
%   dominant eigenvalue decay --- that makes convergence happen.)
%
% \textbf{Trick answer:} B (the trace of $A$ is $1 + 0 = 1$, which 
%   is indeed the sum of eigenvalues. But the trace being $1$ has 
%   no mechanism for ``forcing'' the ratio toward the larger eigenvalue. 
%   The trace constrains the eigenvalue sum; it says nothing about 
%   which eigenvalue dominates the iteration. Convergence is governed 
%   by eigenvalue \emph{magnitudes}, not their sum.)
%
% \textbf{Trick answer:} C ($A = \begin{bmatrix} 1 & 1 \\ 1 & 0 
%   \end{bmatrix}$ is indeed symmetric, so the Spectral Theorem does 
%   apply. However, symmetry is not the reason the dominant eigenvalue 
%   controls the ratio. Any diagonalizable matrix with a strictly 
%   dominant eigenvalue exhibits the same behavior --- symmetry 
%   provides orthogonal eigenvectors and real eigenvalues, but the 
%   convergence mechanism is eigenvalue magnitude separation, which 
%   works identically for non-symmetric matrices. A student selecting 
%   this over-attributes the convergence to a structural property 
%   that is present but irrelevant.)
%
% \textbf{Trick answer:} A (positivity of $\lambda_1$ and positivity 
%   of Fibonacci numbers are both true, but the convergence has 
%   nothing to do with sign matching. If the initial condition had 
%   been chosen so that $c_1 < 0$, the iterates would eventually 
%   become negative, but the \emph{ratio} would still converge to 
%   $\lambda_1$. Moreover, $\lambda_2$ alternates sign when 
%   raised to powers, which does not prevent it from being the 
%   dominant eigenvalue in a different matrix. The sign of the 
%   eigenvalue is not the operative criterion; the magnitude is.)

%%%%%%%%%%%%%%%%%%%%%%%%%%%%%%%%%%%%%%%%%%%%%%%%%%%%
\newpage
%%%%%%%%%%%%%%%%%%%%%%%%%%%%%%%%%%%%%%%%%%%%%%%%%%%%

{\bf PROBLEM 15:}
% WEEK = 8
% TOPICS = Complex eigenvalues, real/imaginary parts, ODE solution behavior, oscillation frequency, growth/decay
%
A $2 \times 2$ real matrix $A$ has complex eigenvalues $\lambda = \alpha \pm i\beta$ 
with $\alpha < 0$ and $\beta > 0$. The solution to $\dfrac{d\mathbf{x}}{dt} = A\mathbf{x}$ 
exhibits oscillations whose amplitude decreases over time.
%
Now consider a modified matrix $A'$ with eigenvalues $\lambda' = \alpha \pm 2i\beta$.
%
Compared to the original system, solutions of $\dfrac{d\mathbf{x}}{dt} = A'\mathbf{x}$ will:

\begin{enumerate}[(A)]
\item Oscillate at twice the frequency and also decay faster, since the eigenvalues are farther from the origin in the complex plane
\item Have identical behavior, since only the real part of the eigenvalue affects the solution of a real ODE
\item Oscillate at the same frequency but decay faster, since the eigenvalues have larger imaginary part
\item Oscillate at twice the frequency and decay at half the rate, since energy is transferred from decay into oscillation
\item Oscillate at twice the frequency with the same rate of amplitude decay
\end{enumerate}

% \textbf{Correct Answer:} E
%
% \textbf{Explanation:} For complex eigenvalues $\alpha \pm i\beta$, the 
% real basis solutions have the form 
% $e^{\alpha t}\cos(\beta t)$ and $e^{\alpha t}\sin(\beta t)$.
% The real part $\alpha$ controls the rate of exponential growth or decay 
% (the envelope), while the imaginary part $\beta$ controls the frequency 
% of oscillation. These two roles are completely independent.
%
% In the modified system, $\alpha$ is unchanged, so the amplitude decay 
% rate is identical. The imaginary part has doubled from $\beta$ to 
% $2\beta$, so the oscillation frequency doubles. The solution form 
% becomes $e^{\alpha t}\cos(2\beta t)$ and $e^{\alpha t}\sin(2\beta t)$:
% same exponential envelope, faster oscillations within it.
%
% The key conceptual point: $\alpha$ and $\beta$ enter the solution in 
% completely separate places --- $\alpha$ in the exponential prefactor, 
% $\beta$ inside the trigonometric functions. Changing one does not 
% affect the other.
%
% \textbf{Partial credit:} A (correctly identifies that doubling the 
%   imaginary part doubles the oscillation frequency --- this is the 
%   harder of the two observations. However, the student then incorrectly 
%   concludes that the decay rate also increases, likely reasoning from 
%   the magnitude $|\lambda'| = \sqrt{\alpha^2 + 4\beta^2} > 
%   \sqrt{\alpha^2 + \beta^2} = |\lambda|$ and conflating eigenvalue 
%   magnitude with decay rate. This reflects correct surface-level 
%   reasoning about frequency but a failure to understand that real 
%   and imaginary parts play independent roles in the solution structure. 
%   A student selecting this has learned that imaginary part $\to$ 
%   frequency, but has not internalized the deeper fact that 
%   $e^{\alpha t}$ and $\cos(\beta t)$ are structurally separate.)
%
% \textbf{Trick answer:} C (reverses which part does what; the imaginary 
%   part controls frequency, not decay rate. A student who selects this 
%   may be confusing the roles of $\alpha$ and $\beta$.)
%
% \textbf{Trick answer:} A (see partial credit note above; the ``farther 
%   from the origin'' reasoning is geometrically intuitive but wrong --- 
%   it's $\alpha$ alone, not $|\lambda|$, that governs decay.)
%
% \textbf{Trick answer:} D (invokes a plausible-sounding but fictitious 
%   ``conservation'' principle --- there is no energy trade-off between 
%   oscillation and decay. The real and imaginary parts act independently 
%   in the solution structure.)
%
% \textbf{Trick answer:} B (confuses the fact that we work with real 
%   basis solutions with the claim that the imaginary part is irrelevant. 
%   The imaginary part is essential: it determines the oscillation 
%   frequency even though the final solution is real-valued.)

\divider

{\bf PROBLEM 16:}
% WEEK = 7
% TOPICS = Coupled systems, diagonalization, decoupling, matrix exponential, eigenbasis
%
Consider the coupled system $\dfrac{d\mathbf{x}}{dt} = A\mathbf{x}$ where 
$A$ is a $2 \times 2$ diagonalizable matrix with distinct real eigenvalues 
$\lambda_1 \neq \lambda_2$ and corresponding eigenvectors 
$\mathbf{v}_1, \mathbf{v}_2$.

We solve this by writing $A = V\Lambda V^{-1}$ and substituting 
$\mathbf{x} = V\mathbf{y}$ to obtain the transformed system 
$\dfrac{d\mathbf{y}}{dt} = \Lambda\mathbf{y}$.
%
What is the essential reason this change of variables helps?

\begin{enumerate}[(A)]
\item It orthogonalizes the eigenvectors, which is necessary for the 
      solution components to evolve independently
\item It makes the matrix $\Lambda$ invertible, which is required to 
      compute the solution $\mathbf{x}(t) = e^{At}\mathbf{x}(0)$
\item It eliminates the eigenvalues, replacing them with simpler constants 
      that are easier to exponentiate
\item Since $\Lambda$ is diagonal, the transformed system consists of 
      independent scalar equations $dy_i/dt = \lambda_i y_i$, each of which 
      can be solved separately
\item It converts the differential equation into an algebraic equation, 
      removing the need to compute any exponentials
\end{enumerate}

% \textbf{Correct Answer:} D
%
% \textbf{Explanation:} The fundamental benefit of diagonalization for 
% solving $d\mathbf{x}/dt = A\mathbf{x}$ is decoupling. In the 
% original coordinates, the equations for the components of $\mathbf{x}$ 
% are intertwined --- each component's derivative depends on the other 
% components. After substituting $\mathbf{x} = V\mathbf{y}$, the 
% system becomes $d\mathbf{y}/dt = \Lambda\mathbf{y}$, which, because 
% $\Lambda$ is diagonal, breaks apart into $n$ independent scalar 
% equations:
%   $$\frac{dy_1}{dt} = \lambda_1 y_1, \quad 
%     \frac{dy_2}{dt} = \lambda_2 y_2$$
% Each has the immediate solution $y_i(t) = y_i(0)e^{\lambda_i t}$. 
% This is something we know how to solve from one-variable calculus. 
% The solution in original coordinates is then recovered via 
% $\mathbf{x}(t) = V\mathbf{y}(t)$.
%
% The conceptual point: diagonalization transforms a coupled system 
% (hard) into independent scalar systems (easy).
%
% \textbf{Partial credit:} A (recognizes that independence of solution 
%   components is the goal, which is correct. However, it incorrectly 
%   claims that orthogonality of eigenvectors is what produces this 
%   independence. In fact, the eigenvectors need not be orthogonal 
%   at all --- any diagonalizable matrix, even one with non-orthogonal 
%   eigenvectors, yields a decoupled system in the eigenbasis. It is 
%   the \emph{diagonality} of $\Lambda$ that decouples the equations, 
%   not any orthogonality property of $V$. A student choosing (E) 
%   understands the goal but conflates diagonalization with 
%   orthogonal diagonalization.)
%
% \textbf{Trick answer:} B ($\Lambda$ may or may not be invertible --- 
%   an eigenvalue could be zero --- and invertibility is not the point. 
%   The matrix exponential $e^{At}$ is always well-defined and 
%   invertible regardless. This answer confuses a property of $\Lambda$ 
%   with the reason diagonalization helps.)
%
% \textbf{Trick answer:} C (the eigenvalues are not eliminated; they 
%   appear directly as the diagonal entries of $\Lambda$ and in the 
%   exponential solutions $e^{\lambda_i t}$. The change of variables 
%   isolates the eigenvalues, it does not remove them.)
%
% \textbf{Trick answer:} E (the problem remains a differential equation 
%   after the substitution --- it just becomes a simpler one. We still 
%   compute exponentials ($e^{\lambda_i t}$), but they are scalar 
%   exponentials rather than matrix exponentials.)
%
% \textbf{Trick answer:} A (see partial credit; the student has the 
%   right intuition about independence but the wrong mechanism)


%%%%%%%%%%%%%%%%%%%%%%%%%%%%%%%%%%%%%%%%%%%%%%%%%%%%
\newpage
%%%%%%%%%%%%%%%%%%%%%%%%%%%%%%%%%%%%%%%%%%%%%%%%%%%%


{\bf PROBLEM 17:}
% WEEK = 8
% TOPICS = Higher-order ODEs, characteristic polynomial, basis solutions, complex eigenvalues, repeated eigenvalues
%
A fourth-order linear homogeneous ODE has characteristic polynomial
$$p(\lambda) = (\lambda - 1)(\lambda + 2)^2(\lambda + 3) = 0$$
%
Which of the following is the correct general solution?

\begin{enumerate}[(A)]
\item $x(t) = c_1 e^{t} + c_2 e^{-2t} + c_3 e^{-3t}$
\item $x(t) = c_1 e^{t} + c_2 e^{-2t} + c_3 e^{-2t}\cos(t) + c_4 e^{-3t}$
\item $x(t) = c_1 e^{t} + c_2 e^{-2t} + c_3 te^{-2t} + c_4 e^{-3t}$
\item $x(t) = c_1 e^{t} + c_2 e^{-2t} + c_3 t^2 e^{-2t} + c_4 e^{-3t}$
\item $x(t) = c_1 e^{t} + c_2 e^{-2t} + c_3 te^{-2t} + c_4 te^{-3t}$
\end{enumerate}

% \textbf{Correct Answer:} C
%
% \textbf{Explanation:} The characteristic polynomial has four roots 
% (counting multiplicity): $\lambda = 1$ (simple), $\lambda = -2$ 
% (repeated with algebraic multiplicity 2), and $\lambda = -3$ (simple).
%
% The rules for constructing basis solutions:
%   - Each simple real root $\lambda$ contributes $e^{\lambda t}$.
%   - A repeated real root $\lambda$ with algebraic multiplicity $k$ 
%     contributes $e^{\lambda t}, \, te^{\lambda t}, \, \ldots, \, 
%     t^{k-1}e^{\lambda t}$.
%   - A complex conjugate pair $\alpha \pm i\beta$ contributes 
%     $e^{\alpha t}\cos(\beta t)$ and $e^{\alpha t}\sin(\beta t)$.
%
% Here, $\lambda = -2$ appears twice, producing both $e^{-2t}$ and 
% $te^{-2t}$. Combined with $e^{t}$ and $e^{-3t}$ from the simple 
% roots, we get exactly four basis solutions --- matching the 
% fourth-order equation.
%
% \textbf{Partial credit:} NONE
%
% \textbf{Trick answer:} A (ignores multiplicity entirely and produces 
%   only 3 basis solutions for a 4th-order equation; this reflects a 
%   more fundamental gap --- the student does not recognize that 
%   algebraic multiplicity demands additional basis solutions)
%
% \textbf{Trick answer:} B (introduces $\cos(t)$ alongside $e^{-2t}$, 
%   as if the repeated real root produces oscillatory behavior. This 
%   conflates repeated real eigenvalues with complex eigenvalues --- 
%   trigonometric terms arise only from imaginary parts, and 
%   $\lambda = -2$ is purely real.)
%
% \textbf{Trick answer:} D (the student correctly identifies that: 
%   (1) a fourth-order equation needs four basis solutions, (2) the 
%   repeated root $\lambda = -2$ produces a polynomial-exponential 
%   term, and (3) this term is attached specifically to $\lambda = -2$ 
%   and not to the simple roots. The only error is the power of $t$: 
%   for multiplicity $k$, the prefactors run $t^0, t^1, \ldots, t^{k-1}$, 
%   so multiplicity 2 gives $e^{-2t}$ and $te^{-2t}$, not $t^2e^{-2t}$. 
%   This is a minor indexing error reflecting genuine but imperfect 
%   understanding of the repeated-root mechanism.)
%
% \textbf{Trick answer:} E (correctly produces $te^{-2t}$ for the 
%   repeated root, but also introduces a spurious $te^{-3t}$ for the 
%   simple root $\lambda = -3$. This suggests the student knows 
%   that polynomial prefactors appear somewhere but applies them 
%   indiscriminately rather than only to repeated eigenvalues.)


\divider


{\bf PROBLEM 18:}
% WEEK = 7
% TOPICS = Eigenvalue zero, singularity, invertibility, ker(A), matrix exponential
%
A $4 \times 4$ matrix $A$ has characteristic polynomial $p(\lambda) = \lambda^2(\lambda - 3)(\lambda + 1)$.
%
Which of the following must be true?
%
\begin{enumerate}[(A)]
\item $\ker(A) \neq \{\mathbf{0}\}$
\item $A^{-1}$ exists
\item $\text{tr}(A) = 0$
\item $e^{A}$ is singular
\item $A$ is diagonalizable
\end{enumerate}

% \textbf{Correct Answer:} A
% \textbf{Explanation:} The characteristic polynomial $\lambda^2(\lambda-3)(\lambda+1)$ has roots $\lambda = 0$ (algebraic multiplicity 2), $\lambda = 3$, and $\lambda = -1$. Since $0$ is an eigenvalue, there exists a nonzero $\mathbf{v}$ with $A\mathbf{v} = \mathbf{0}$, so $\ker(A)$ is nontrivial. Equivalently, $\det(A) = 0 \cdot 0 \cdot 3 \cdot (-1) = 0$, so $A$ is singular.
% \textbf{Partial credit:} None
% \textbf{Wrong:} E — The repeated eigenvalue $0$ has algebraic multiplicity 2, but the geometric multiplicity could be 1 or 2. If geometric multiplicity is 1, $A$ is not diagonalizable (defective). We cannot determine this from the characteristic polynomial alone.
% \textbf{Wrong:} B — Since $\det(A) = 0$, $A^{-1}$ does NOT exist. This is the opposite of what's true.
% \textbf{Trick answer:} C — The trace equals the sum of eigenvalues: $\text{tr}(A) = 0 + 0 + 3 + (-1) = 2 \neq 0$. A student who confuses trace (sum) with determinant (product) might select this.
% \textbf{Trick answer:} D — $e^A$ is ALWAYS nonsingular (since $\det(e^A) = e^{\text{tr}(A)} = e^2 \neq 0$). This is a common misconception: students think singularity of $A$ implies singularity of $e^A$.

\divider

{\bf PROBLEM 19:}
% WEEK = 7
% TOPICS = Matrix powers via diagonalization, eigendecomposition of vectors, 
%          eigenvalue powers, sign of (-1)^k for even k
%
A diagonalizable $2 \times 2$ matrix $A$ has eigenvalues $\lambda_1 = 2$ 
and $\lambda_2 = -1$ with corresponding eigenvectors $\mathbf{v}_1$ 
and $\mathbf{v}_2$.

Given an initial vector $\mathbf{x}_0 = 3\mathbf{v}_1 + 5\mathbf{v}_2$, 
what is $A^{100}\mathbf{x}_0$?

\begin{enumerate}[(A)]
\item $(3 \cdot 2^{100} + 5)\,\mathbf{v}_1$
\item $2^{100}(3\,\mathbf{v}_1 + 5\,\mathbf{v}_2)$
\item $3 \cdot 2^{100}\,\mathbf{v}_1 - 5\,\mathbf{v}_2$
\item $3 \cdot 2^{100}\,\mathbf{v}_1$
\item None of the above
\end{enumerate}

% \textbf{Correct Answer:} E
%
% \textbf{Explanation:} The diagonalization $A = V\Lambda V^{-1}$ gives 
% $A^k = V\Lambda^k V^{-1}$. When we expand $\mathbf{x}_0$ in the 
% eigenbasis, the matrix power acts on each component independently:
%   $$A^{100}\mathbf{x}_0 = 3\lambda_1^{100}\,\mathbf{v}_1 + 
%     5\lambda_2^{100}\,\mathbf{v}_2 = 3 \cdot 2^{100}\,\mathbf{v}_1 + 
%     5 \cdot (-1)^{100}\,\mathbf{v}_2$$
%
% Since $(-1)^{100} = +1$ (even power), this simplifies to:
%   $$A^{100}\mathbf{x}_0 = 3 \cdot 2^{100}\,\mathbf{v}_1 + 
%     5\,\mathbf{v}_2$$
%
% This expression does not appear among choices (A)--(D), so the 
% answer is (E).
%
% The key ideas tested:
%   (1) $A^k$ acts on each eigenvector component by raising the 
%       corresponding eigenvalue to the $k$-th power: the 
%       coefficient $c_i$ becomes $c_i \lambda_i^k$.
%   (2) $(-1)^{100} = +1$, not $-1$. Even powers of $-1$ are 
%       positive.
%   (3) The two eigenvector components evolve independently --- 
%       they do not merge, cancel, or interact.
%
% \textbf{Partial credit:} C (correctly applies the eigendecomposition 
%   formula to get $3 \cdot 2^{100}\,\mathbf{v}_1$ for the first term, 
%   and correctly recognizes that $(-1)^{100}$ has magnitude $1$. The 
%   only error is the sign: $(-1)^{100} = +1$, not $-1$. This is a 
%   minor arithmetic slip reflecting solid conceptual understanding 
%   of how matrix powers act through the eigenbasis, with only the 
%   even/odd distinction missed.)
%
% \textbf{Trick answer:} C (see partial credit; miscomputes 
%   $(-1)^{100} = -1$ instead of $+1$. This is the most common 
%   error: the student correctly sets up the eigendecomposition 
%   but defaults to thinking $(-1)^k$ is always $-1$, forgetting 
%   to check the parity of $k$.)
%
% \textbf{Trick answer:} B (treats the dominant eigenvalue $2^{100}$ 
%   as a uniform scaling factor applied to the entire vector. This 
%   misses the fundamental point that each eigenvector component 
%   is scaled by its \emph{own} eigenvalue power. The matrix does 
%   not act like a scalar multiple --- different directions are 
%   stretched differently.)
%
% \textbf{Trick answer:} D (assumes the $\mathbf{v}_2$ component 
%   vanishes because $(-1)^{100}$ is ``small'' or ``decays.'' In 
%   fact, $|{-1}|^{100} = 1$: the magnitude of the $\mathbf{v}_2$ 
%   component is perfectly preserved. This distractor catches 
%   students who confuse $|\lambda| < 1$ (which causes decay) with 
%   $\lambda < 0$ (which causes sign alternation but no decay). 
%   The eigenvalue $-1$ has magnitude exactly $1$, so it neither 
%   grows nor decays.)
%
% \textbf{Trick answer:} A (collapses both eigenvector components 
%   into the dominant direction $\mathbf{v}_1$, as if the large 
%   factor $2^{100}$ somehow absorbs the $\mathbf{v}_2$ component. 
%   This reflects a misunderstanding of eigendecomposition: the 
%   components along different eigenvectors are independent and 
%   never ``merge'' into a single direction, no matter how large 
%   one eigenvalue power becomes relative to the other.)


%%%%%%%%%%%%%%%%%%%%%%%%%%%%%%%%%%%%%%%%%%%%%%%%%%%%
\newpage
%%%%%%%%%%%%%%%%%%%%%%%%%%%%%%%%%%%%%%%%%%%%%%%%%%%%


{\bf PROBLEM 20:}
% WEEK = 8
% TOPICS = QR algorithm, similarity transformation, eigenvalue preservation
% NOTE: Complements P15 (which asks WHAT the QR iterates converge to);
%       this asks WHY the eigenvalues are preserved at each step.
%
The QR algorithm computes $A_{k+1} = R_k Q_k$ from the QR factorization 
$A_k = Q_k R_k$.
%
A key property is that every iterate $A_{k+1}$ has the same eigenvalues 
as $A_k$.
%
Why is this true?

\begin{enumerate}[(A)]
\item Because $Q_k$ is orthogonal, and orthogonal transformations 
      preserve all matrix properties including eigenvalues
\item Because the QR factorization itself preserves eigenvalues: 
      $Q_k$ and $R_k$ each have the same eigenvalues as $A_k$
\item Because $A_{k+1} = Q_k^{-1} A_k Q_k$, which is a similarity 
      transformation
\item Because the diagonal entries of $R_k$ are the eigenvalues of 
      $A_k$, and these entries are reused to construct $A_{k+1}$
\item Because $R_k Q_k$ and $Q_k R_k$ are transposes of each other, 
      and a matrix and its transpose share eigenvalues
\end{enumerate}

% \textbf{Correct Answer:} C
%
% \textbf{Explanation:} The crucial algebraic identity:
%   $$A_{k+1} = R_k Q_k = (Q_k^{-1} Q_k) R_k Q_k 
%     = Q_k^{-1} (Q_k R_k) Q_k = Q_k^{-1} A_k Q_k$$
% This shows $A_{k+1}$ is obtained from $A_k$ by a similarity 
% transformation (conjugation by $Q_k$). Similar matrices have the 
% same characteristic polynomial and therefore the same eigenvalues. 
% This is the foundational reason the QR algorithm works: it 
% iteratively transforms $A$ toward a simpler form (triangular) 
% while preserving its eigenvalues at every step.
%
% \textbf{Trick answer:} B --- The QR factorization does NOT 
%   individually preserve eigenvalues. The eigenvalues of $Q_k$ 
%   all have modulus $1$ (since $Q_k$ is orthogonal), and the 
%   eigenvalues of $R_k$ are its diagonal entries (since $R_k$ is 
%   upper triangular). Neither set of eigenvalues matches those of 
%   $A_k$ in general. The eigenvalue preservation comes from the 
%   \emph{re-multiplication} $R_kQ_k$, not from the factorization 
%   itself.
%
% \textbf{Trick answer:} A --- ``Orthogonal transformations preserve 
%   all matrix properties'' is far too broad. Orthogonal 
%   transformations preserve eigenvalues (via similarity), inner 
%   products, and norms, but they do \emph{not} preserve all matrix 
%   properties (e.g., they generally change the entries, the sparsity 
%   pattern, and so on). More importantly, the reason is specifically 
%   that $Q_k^{-1}A_kQ_k$ is a \emph{similarity transformation}, 
%   not just that $Q_k$ is orthogonal. Any invertible $P$ gives a 
%   similarity $P^{-1}A_kP$ preserving eigenvalues; orthogonality 
%   is helpful for numerical stability but is not the reason 
%   eigenvalues are preserved.
%
% \textbf{Trick answer:} D --- The diagonal entries of $R_k$ are 
%   related to norms from the Gram-Schmidt process, \emph{not} to 
%   the eigenvalues of $A_k$. This conflates two different roles of 
%   diagonal entries: in a triangular matrix obtained by 
%   \emph{similarity} (like the Schur form), diagonal entries are 
%   eigenvalues. In the $R$ factor from QR \emph{factorization}, 
%   they are not.
%
% \textbf{Trick answer:} E --- $R_kQ_k$ and $Q_kR_k$ are generally 
%   \emph{not} transposes of each other. While it is true that a 
%   matrix and its transpose share eigenvalues (same characteristic 
%   polynomial), this fact is irrelevant here because the transpose 
%   relationship does not hold. This conflates two unrelated 
%   algebraic facts.
%
% \textbf{Partial credit:} None

\divider


{\bf PROBLEM 21:}
% WEEK = 7
% TOPICS = Eigenvalues, trace, determinant, characteristic polynomial, properties of eigenvalues
%
A $3 \times 3$ real matrix $A$ satisfies $\operatorname{tr}(A) = 7$, $\det(A) = 0$, and is known to have an eigenvalue $\lambda_1 = 3$.
%
Which of the following must be the remaining eigenvalues?

\begin{enumerate}[(A)]
\item $\lambda_2 = 4, \; \lambda_3 = 0$
\item $\lambda_2 = 4, \; \lambda_3 = 3$
\item $\lambda_2 = 2+i, \; \lambda_3 = 2-i$
\item $\lambda_2 = 0, \; \lambda_3 = 0$
\item $\lambda_2 = 7, \; \lambda_3 = 0$
\end{enumerate}

% \textbf{Correct Answer:} A
%
% \textbf{Explanation:} Two constraints must be satisfied simultaneously:
%   (1) $\operatorname{tr}(A) = \lambda_1 + \lambda_2 + \lambda_3 = 7$, 
%       so $\lambda_2 + \lambda_3 = 4$.
%   (2) $\det(A) = \lambda_1 \cdot \lambda_2 \cdot \lambda_3 = 0$, 
%       so at least one of $\lambda_2, \lambda_3$ must be zero.
%   Additionally, since $A$ is a real matrix, complex eigenvalues must 
%   come in conjugate pairs.
%
%   Checking each choice:
%   (A) $4 + 0 = 4$ and $3 \cdot 4 \cdot 0 = 0$. 
%       Both conditions satisfied. VALID.
%   (B) $7 + 0 = 7 \neq 4$. Fails the trace condition.
%   (C) $(2+i) + (2-i) = 4$ and conjugate pair is valid for a real matrix. 
%       But $\det = 3(2+i)(2-i) = 3 \cdot 5 = 15 \neq 0$. 
%       Fails the determinant condition.
%   (D) $0 + 0 = 0 \neq 4$. Fails the trace condition.
%   (E) $4 + 3 = 7 \neq 4$ (i.e., $3 + 4 + 3 = 10 \neq 7$). 
%       Fails the trace condition.
%
% \textbf{Partial credit:} C (student correctly verified the trace condition 
%   AND recognized the conjugate pair requirement for real matrices, but 
%   failed to check the determinant condition --- reflects partial 
%   understanding: two out of three constraints checked.)
%
% \textbf{Trick answer:} C (satisfies trace and conjugate-pair requirements 
%   but violates $\det = 0$; catches students who check only one or two conditions)
% \textbf{Trick answer:} E (has a zero eigenvalue, appealing if focused only 
%   on $\det = 0$, but violates the trace)
% \textbf{Trick answer:} B (plausible-looking integers, but $3+4+3=10\neq 7$; 
%   catches careless arithmetic)
% \textbf{Trick answer:} D (double zero is tempting for students who over-focus 
%   on $\det = 0$ without checking the sum)
 

\divider


{\bf PROBLEM 22:}
% WEEK = 7
% TOPICS = Diagonalizability, algebraic vs geometric multiplicity, 
%          sufficient conditions, repeated eigenvalues
%
A $3 \times 3$ matrix $A$ has eigenvalues $\lambda_1 = 4$, 
$\lambda_2 = 4$, $\lambda_3 = 1$.
%
Which additional piece of information would guarantee that $A$ is 
diagonalizable?
\begin{enumerate}[(A)]
\item $\ker(A - 4I)$ is one-dimensional
\item $A$ is upper triangular
\item The columns of $A$ are linearly independent
\item The eigenspace for $\lambda = 4$ is two-dimensional
\item None of the above: a repeated eigenvalue means $A$ cannot 
      be diagonalized
\end{enumerate}
% \textbf{Correct Answer:} D
%
% \textbf{Explanation:} The eigenvalue $\lambda = 4$ has algebraic 
% multiplicity $2$, so diagonalizability hinges on whether its 
% geometric multiplicity also equals $2$. If the eigenspace for 
% $\lambda = 4$ is two-dimensional, then geometric multiplicity 
% matches algebraic multiplicity for this eigenvalue. The eigenvalue 
% $\lambda = 1$ has algebraic multiplicity $1$, so its geometric 
% multiplicity is automatically $1$. With geometric $=$ algebraic 
% for \emph{every} eigenvalue, $A$ is diagonalizable.
%
% \textbf{Partial credit:} NONE 

% \textbf{Trick answer:} A --- correctly identifies 
%   $\ker(A - 4I)$ as the eigenspace for the repeated eigenvalue 
%   and recognizes that this eigenspace is the object that determines 
%   diagonalizability. However, a one-dimensional eigenspace means 
%   geometric multiplicity $1 < 2 =$ algebraic multiplicity, which 
%   actually guarantees $A$ is \emph{not} diagonalizable. The student 
%   has the right framework but the wrong dimension threshold.)
%
% \textbf{Trick answer:} B --- Upper triangular matrices display 
%   their eigenvalues on the diagonal, but this gives no information 
%   about eigenspace dimensions. A Jordan block such as 
%   $\begin{bmatrix} 4 & 1 \\ 0 & 4 \end{bmatrix}$ is upper 
%   triangular with repeated eigenvalue $4$ yet is not 
%   diagonalizable. Upper triangularity is about the form of the 
%   matrix, not the structure of its eigenspaces.
%
% \textbf{Trick answer:} C --- Since no eigenvalue is zero, 
%   $\det(A) = 4 \cdot 4 \cdot 1 = 16 \neq 0$, so $A$ is already 
%   invertible (i.e., its columns are already linearly independent) 
%   from the given information alone. This tells us nothing new. 
%   Moreover, invertibility is unrelated to diagonalizability: a 
%   matrix can be invertible but not diagonalizable (e.g., a Jordan 
%   block with $\lambda \neq 0$).
%
% \textbf{Trick answer:} E --- This is a common misconception. 
%   Repeated eigenvalues do \emph{not} prevent diagonalizability; 
%   they merely create a \emph{condition} that must be checked. 
%   For instance, $4I_{3\times 3}$ has eigenvalue $4$ with algebraic 
%   multiplicity $3$ and is trivially diagonalizable (it is already 
%   diagonal). Diagonalizability fails only when an eigenspace is 
%   too small, not when an eigenvalue is repeated.


\end{document}